\documentclass[10pt,letterpaper]{article}

% ── Packages ──────────────────────────────────────────────────────────────
\usepackage[utf8]{inputenc}
\usepackage[T1]{fontenc}
\usepackage{lmodern}
\usepackage[margin=1in]{geometry}
\usepackage{amsmath,amssymb,amsfonts}
\usepackage{booktabs}
\usepackage{graphicx}
\usepackage{xcolor}
\usepackage{hyperref}
\usepackage{cleveref}
\usepackage{algorithm}
\usepackage{algpseudocode}
\usepackage{multirow}
\usepackage{enumitem}
\usepackage{caption}
\usepackage{subcaption}
\usepackage{microtype}

% ── Hyperref setup ────────────────────────────────────────────────────────
\hypersetup{
  colorlinks=true,
  linkcolor=blue!70!black,
  citecolor=blue!70!black,
  urlcolor=blue!70!black,
  pdftitle={Fused Block-Sparse Attention with Learned Content-Dependent Selection},
  pdfauthor={Miroslav Drbal}
}

% ── Custom commands ───────────────────────────────────────────────────────
\newcommand{\Ocomplex}{\mathcal{O}}
\newcommand{\R}{\mathbb{R}}
\newcommand{\topk}{\operatorname{top\text{-}k}}
\newcommand{\softmax}{\operatorname{softmax}}
\newcommand{\LogSigmoid}{\operatorname{log\text{-}sigmoid}}
\newcommand{\InfoNCE}{\textsc{InfoNCE}}

% ── Title block ───────────────────────────────────────────────────────────
\title{Fused Block-Sparse Attention with Learned\\
Content-Dependent Selection:\\[4pt]
\large A Triton Kernel for Memory-Efficient Long-Context Inference}

\author{%
  Miroslav Drbal\\
  \texttt{mdrbal@nymfe.net}\\
  Gen Digital, Inc.~\textbar~\texttt{miroslav.drbal@gendigital.com}
}
\date{}

\begin{document}
\maketitle

% ─═ Abstract ══════════════════════════════════════════════════════════════
\begin{abstract}
We present a block-sparse attention mechanism with a fused Triton kernel that
achieves memory parity with dense FlashAttention while delivering up to
$7.4\times$ latency speedup at 262k-token context length. The block selector
is trained jointly with the model via MoE-style gate coupling and an auxiliary
routing loss, reaching near-perfect block-retrieval accuracy (hit $= 0.99$)
without a separate pre-training stage. On a 45M-parameter language model
trained on TinyStories, the sparse model achieves validation perplexity
15.0 versus 14.4 for the dense baseline, a 4\% gap that narrows to parity at
shorter context windows. We also investigate centroid-based linear selection
($\Ocomplex(N \cdot C)$) as a replacement for the $\Ocomplex(N^2/B_s^2)$
block-pair selection, finding that InfoNCE contrastive training fixes a
critical collapse failure in the routing loss (hit $0.19 \to 1.0$) but that
task accuracy remains bottlenecked by selection precision, not routing recall.
We further isolate selection from the attention read up to 16M tokens and show
that single-level block-pair selection has an $\Ocomplex(N^{4/3})$ compute
floor: no block-size schedule yields end-to-end linearity, as the quadratic
only moves between the selection and read stages.
A second, lossless linear selector, LSH bucketing, matches block-sparse accuracy
under supervised routing (0.92 vs.\ 0.96 on MQAR) and far exceeds centroid, but
its unsupervised (gate-only) routing is inert in a 45M language model, matching a
no-routing baseline, leaving reliable unsupervised linear selection open.
All experiments run on a single NVIDIA GB10 Grace-Blackwell GPU.
\end{abstract}

% ─═ 1. Introduction ═══════════════════════════════════════════════════════
\section{Introduction}\label{sec:intro}

Transformer-based language models face a quadratic attention bottleneck:
standard softmax attention scales as $\Ocomplex(N^2)$ in sequence length $N$,
limiting practical context windows. While FlashAttention
\cite{dao2022flashattention,dao2023flashattention2} reduces the \emph{constant}
factor by tiling computations in SRAM, the asymptotic cost remains quadratic.

Content-dependent sparse attention offers a path to genuine sub-quadratic
scaling by selecting a small subset of key/value blocks per query block.
Mechanisms such as Native Sparse Attention (NSA)~\cite{yuan2025nsa}, MoBA
\cite{lu2025moba}, and the Routing Transformer~\cite{roy2021routing} learn to
route queries to relevant key blocks, reducing the attention cost to
$\Ocomplex(N \cdot k \cdot B_s)$ where $k$ is the number of selected blocks
and $B_s$ is the block size.

However, two practical barriers remain:

\begin{enumerate}[leftmargin=*,itemsep=2pt]
\item \textbf{Selector training.} Discrete top-$k$ selection severs the
  gradient path to the selector parameters. Without auxiliary supervision, the
  selector remains a frozen random projection.
\item \textbf{Gather overhead.} Even with correct selection, the reference
  PyTorch implementation materializes gathered key/value tensors, consuming up
  to $40\times$ more memory than dense FlashAttention and negating the
  theoretical savings.
\end{enumerate}

We address both barriers. For~(1), we use MoE-style gate coupling
\cite{shazeer2017moe} -adding $\LogSigmoid(\text{score})$ as a bias to
attention logits, so the selector receives task gradients. An auxiliary
routing loss (cross-entropy on the known needle block for synthetic tasks,
InfoNCE contrastive loss~\cite{oord2018infonce} for centroid routing) provides
cold-start supervision. For~(2), we implement a fused Triton kernel
\cite{tillet2019triton} that reads selected block indices, gathers key/value
blocks one at a time from the KV cache, and computes online-softmax in
SRAM, never materializing the gathered tensor.

\paragraph{Contributions.}
\begin{itemize}[leftmargin=*,itemsep=2pt]
\item A \textbf{fused Triton kernel} for block-sparse attention (forward +
  backward) achieving memory parity with dense FlashAttention
  ($1.11$~GB vs.\ $1.12$~GB at 262k tokens) and $7.4\times$ latency speedup
  (\Cref{sec:triton}).
\item \textbf{Gate-coupled selector training} with InfoNCE routing loss,
  fixing a collapse failure in centroid-based linear selection
  (hit $0.19 \to 1.0$) and characterizing the precision/recall tradeoff
  (\Cref{sec:centroid}).
\item An \textbf{end-to-end language model} (45M params, TinyStories
  \cite{eldan2023tinystories}) trained with the fused kernel, achieving val
  perplexity 15.0 vs.\ 14.4 dense (\Cref{sec:lm}).
\item A \textbf{negative result}: linear-selection centroid routing achieves
  perfect retrieval recall but cannot match block-sparse task accuracy,
  identifying representation quality under coarse routing as the bottleneck
  (\Cref{sec:centroid}).
\item A \textbf{scaling analysis} isolating selection from the attention read
  to 16M tokens: the read is linear, but single-level block-pair selection has
  an $\Ocomplex(N^{4/3})$ floor that no block-size schedule can beat and that
  exhausts memory before any latency crossover (\Cref{sec:floor}).
\end{itemize}

% ─═ 2. Background ═════════════════════════════════════════════════════════
\section{Background and Related Work}\label{sec:related}

\paragraph{Sparse attention.}
Fixed-pattern sparsity (Longformer~\cite{beltagy2020longformer}, BigBird
\cite{zaheer2020bigbird}) restricts attention to local windows plus global
tokens but cannot solve content-dependent retrieval tasks like multi-query
associative recall (MQAR)~\cite{arora2024mqar,olsson2022induction}.
Content-dependent methods -Routing Transformer~\cite{roy2021routing},
Reformer~\cite{kitaev2020reformer}, Sparse Sinkhorn Attention
\cite{tay2020sinkhorn}, NSA~\cite{yuan2025nsa}, MoBA~\cite{lu2025moba} -learn
to select relevant blocks but face the selector-gradient problem.

\paragraph{Inference-time sparse selection.}
A parallel line keeps a dense-trained model and sparsifies only at inference,
selecting a query-dependent subset of the KV cache: Quest~\cite{tang2024quest}
(top-$k$ pages), H2O~\cite{zhang2023h2o} (heavy-hitter eviction), and
RetrievalAttention~\cite{liu2024retrievalattention} (attention-aware ANN over
keys, which reports that plain LSH/ANN underperforms because attention queries
are out-of-distribution). LSH-E~\cite{rabbani2024lshe} uses SimHash for KV
eviction, and DASH-KV~\cite{guo2026dashkv} replaces LSH with a \emph{learned}
asymmetric hash, reporting higher retrieval recall than parameter-free LSH. Our
train-dense / infer-LSH study (\Cref{sec:lsh}) sits in this regime, at block
rather than token granularity and at small scale.

\paragraph{Selector training.}
NSA trains the selector by distilling dense per-block attention mass.
MoBA uses a gating mechanism. Sparse Sinkhorn Attention
\cite{tay2020sinkhorn} uses differentiable sorting via the Sinkhorn operator
to learn latent permutations. We build on the MoE gate-coupling approach
\cite{shazeer2017moe}: the selector score is injected as a log-sigmoid bias
into the attention logits, creating a differentiable path. An auxiliary
routing loss provides cold-start supervision.

\paragraph{Flash-style fusion.}
FlashAttention~\cite{dao2022flashattention} tiles Q/K/V into blocks and
computes online-softmax in SRAM, achieving exact attention with reduced memory
I/O. FlashAttention-2~\cite{dao2023flashattention2} improves parallelism and
work partitioning. Our Triton kernel applies the same dataflow but iterates
over \emph{selected} key blocks via an index list rather than all blocks in
sequence order.

\paragraph{Centroid routing.}
Routing Transformer~\cite{roy2021routing} assigns tokens to $k$-means
clusters. Reformer~\cite{kitaev2020reformer} uses locality-sensitive hashing.
Our CentroidSSA routes blocks through learned centroids
($\Ocomplex(N \cdot C)$) but uses InfoNCE contrastive training
\cite{oord2018infonce} instead of $k$-means, addressing the
non-differentiability of hard cluster assignment.

\paragraph{Triton.}
Triton~\cite{tillet2019triton} is an intermediate language and compiler for
tiled neural network computations, enabling custom GPU kernels at near-CUDA
performance from Python. Our forward and backward kernels are implemented
in Triton 3.7.

% ─═ 3. Method ═════════════════════════════════════════════════════════════
\section{Method}\label{sec:method}

\subsection{Block-Sparse Attention with Gate Coupling}

Given input $x \in \R^{B \times L \times D}$, we partition the sequence into
$n_{\text{blk}} = L / B_s$ blocks of size $B_s$. Per layer:

\begin{enumerate}[leftmargin=*,itemsep=1pt]
\item \textbf{Projections.} $Q, K, V = \text{QKV}(x)$, reshaped to
  $(B, H, L, d_h)$.
\item \textbf{Block scores.} A small MLP selector computes block-level
  summaries:
  \begin{align}
    s_q &= \text{mean}(\text{sel\_q}(x)_{\text{blk}}) \in \R^{B \times n_{\text{blk}} \times d_s} \\
    s_k &= \max\bigl(\text{mean}(\text{sel\_k}(x)_{\text{blk}}),\;
               \max(\text{sel\_k}(x)_{\text{blk}})\bigr) \\
    S &= \max(s_q s_{k,\text{mean}}^\top,\; s_q s_{k,\text{max}}^\top)
      \in \R^{B \times n_{\text{blk}} \times n_{\text{blk}}}
  \end{align}
\item \textbf{Selection.} For each query block $i$, select the top-$k$ key
  blocks by $S[i, \cdot]$, plus the own (diagonal) block. Causal masking
  excludes future blocks. Result:
  $\text{sel} \in \mathbb{Z}^{B \times n_{\text{blk}} \times k_{\text{tot}}}$.
\item \textbf{Gate coupling.} Add $\LogSigmoid(S[i, j])$ as a bias to the
  per-token attention logits for selected block $j$:
  \begin{equation}
    \text{logit}_{i,j} = \frac{Q_i K_j^\top}{\sqrt{d_h}} + \LogSigmoid(S_{i,j})
  \end{equation}
  This makes the output depend differentiably on $S$, so $\text{sel\_q}$ and
  $\text{sel\_k}$ receive task gradients.
\item \textbf{Attend.} Exact scaled-dot-product attention over the selected
  blocks only ($k_{\text{tot}} \cdot B_s$ keys per query block).
\end{enumerate}

\subsection{Auxiliary Routing Loss}

For synthetic tasks (MQAR) where needle positions are known, we supervise the
selector with cross-entropy to the true needle block:
\begin{equation}
  \mathcal{L}_{\text{route}} = \text{CE}\bigl(\softmax(S_{\text{last}}),\;
  \lfloor \text{needle\_pos} / B_s \rfloor\bigr)
\end{equation}
For centroid routing (\Cref{sec:centroid}), we replace this with an InfoNCE
contrastive loss. The total loss is
$\mathcal{L} = \mathcal{L}_{\text{task}} + \lambda_r \mathcal{L}_{\text{route}}$,
with $\lambda_r$ annealed to zero over the second half of training.

\subsection{Complexity}

\begin{itemize}[leftmargin=*,itemsep=1pt]
\item \textbf{Attention:} $\Ocomplex(N \cdot k \cdot B_s)$ -linear in $N$.
\item \textbf{Selection (block-pair):} $\Ocomplex(n_{\text{blk}}^2) =
  \Ocomplex((N/B_s)^2)$, quadratic on the reduced axis, shrunk by $B_s^2$.
\item \textbf{Selection (centroid):} $\Ocomplex(n_{\text{blk}} \cdot C)$, linear.
\item \textbf{Memory (fused kernel):} $\Ocomplex(N)$ -only base Q/K/V,
  no gathered tensor.
\end{itemize}

% ─═ 4. Fused Triton Kernel ════════════════════════════════════════════════
\section{Fused Triton Kernel}\label{sec:triton}

\subsection{Forward Kernel}

Each Triton program handles one $(B, H, q_{\text{blk}})$ tile. It loads the
$B_s \times d_h$ query tile, then iterates over the $k_{\text{tot}}$ selected
key blocks via the index list $\text{sel}$. For each selected block:

\begin{enumerate}[leftmargin=*,itemsep=1pt]
\item Load $K, V$ tiles ($B_s \times d_h$) from the KV cache at the indexed
  block offset.
\item Compute $B_s \times B_s$ attention scores with scale.
\item Apply causal mask (triu) for the own block; mask future blocks entirely.
\item Add gate bias $\LogSigmoid(S_{i,j})$ if enabled.
\item Update online-softmax state (running max $m$, running sum $\ell$,
  accumulator $\text{acc}$) via the standard FlashAttention recurrence
  \cite{dao2022flashattention}.
\end{enumerate}

The output is $\text{acc} / \ell$. No gathered K/V tensor or score matrix is
ever materialized, all intermediate state lives in SRAM.

\subsection{Backward Kernel}

The backward pass follows the FlashAttention-2 recompute pattern
\cite{dao2023flashattention2} with three iterations over the selected blocks:

\begin{enumerate}[leftmargin=*,itemsep=1pt]
\item \textbf{Pass 1:} Recompute forward to obtain final $m$ and $\ell$.
\item \textbf{Pass 2:} Compute $D = \sum_j \text{rowsum}(P_j \cdot dP_j)$,
  the total diagonal term for the softmax Jacobian across all selected blocks.
\item \textbf{Pass 3:} Compute gradients:
  \begin{align}
    dV &= P^\top dO \\
    dP &= dO \, V^\top \\
    dS &= P \odot (dP - D) \cdot \text{scale} \\
    dQ &\mathrel{+}= dS \, K, \quad dK \mathrel{+}= dS^\top Q
  \end{align}
  $dK$ and $dV$ are written via \texttt{atomic\_add} (multiple query blocks
  may select the same key block); $dQ$ is written directly.
\end{enumerate}

\subsection{Integration}

A custom \texttt{torch.autograd.Function} (\texttt{BSAttnFunction}) routes
forward and backward through the Triton kernels. The
\texttt{BlockSparseSSA(use\_triton=True)} flag switches between the PyTorch
reference and the fused path. Non-contiguous tensors from the QKV projection
are made contiguous before the kernel call.

\subsection{Verification}

\begin{itemize}[leftmargin=*,itemsep=1pt]
\item \textbf{Forward:} $9.77 \times 10^{-4}$ max diff vs.\ PyTorch SDPA
  (fp16, $B{=}2, H{=}4, L{=}64, d_h{=}32$, $k{=}3$, causal).
\item \textbf{Backward:} $dQ$: $4.0 \times 10^{-3}$, $dK$: $4.4 \times 10^{-3}$,
  $dV$: $3.7 \times 10^{-3}$ vs.\ PyTorch autograd (fp32, $k{=}3$, causal).
\item \textbf{Integration:} forward $4.5 \times 10^{-4}$, gradient
  $1.1 \times 10^{-3}$ vs.\ PyTorch path (same model, toggle
  \texttt{use\_triton}).
\item \textbf{Gate bias:} $9.77 \times 10^{-4}$ with gate enabled.
\end{itemize}

% ─═ 5. Centroid Routing ═══════════════════════════════════════════════════
\section{Linear Selection: Centroid and LSH}\label{sec:centroid}

Block-pair selection scales as $\Ocomplex(n_{\text{blk}}^2)$. To remove this
quadratic term, we propose CentroidSSA: route blocks through $C$ learned
centroids, analogous to $k$-means routing in the Routing Transformer
\cite{roy2021routing} but with differentiable training.

\subsection{Mechanism}

Each block is assigned to its nearest centroid (argmax of
$s_k \cdot \text{cent}^\top$). Each query block routes to its top-$c_{pq}$
centroids and attends to all blocks in those buckets (capacity $c_{\text{cap}}$
per bucket). Selection cost: $\Ocomplex(n_{\text{blk}} \cdot C)$, linear.

\subsection{The Circular Loss Failure}

The original routing-alignment loss was:
\begin{equation}
  \text{consensus} = \operatorname{argmax}(s_q + s_k), \quad
  \mathcal{L} = \text{CE}(s_q, \text{consensus}) + \text{CE}(s_k, \text{consensus})
\end{equation}
This is \emph{circular}: at cold start, $s_q$ and $s_k$ are near-uniform, so
$\operatorname{argmax}$ is arbitrary noise. The loss collapsed to $\ln C$ (uniform), and
retrieval hit stayed at 0.19.

\subsection{InfoNCE Fix}

We replace the circular loss with an InfoNCE contrastive objective
\cite{oord2018infonce}:
\begin{equation}
  \mathcal{L}_{\text{InfoNCE}} = -\log \frac{
    \exp(\text{sim}(q_{\text{last}}, k_{\text{needle}}) / \tau)
  }{
    \sum_{j} \exp(\text{sim}(q_{\text{last}}, k_j) / \tau)
  }
\end{equation}
where $\text{sim}(q, k) = \softmax(q/\tau) \cdot \softmax(k/\tau)^\top$ and
$\tau = 0.5$. The positive is the needle key block; all other key blocks are
negatives. This is fully differentiable and has no argmax dependency.

\subsection{Capacity Sweep}

\begin{table}[h]
\centering
\caption{Centroid capacity sweep (8-pair MQAR, 10k steps, curriculum).
$C$ = centroids, cap = bucket capacity.}
\label{tab:centroid}
\small
\begin{tabular}{rrll}
\toprule
$C$ & cap & hit & acc \\
\midrule
16 & 4 & 0.71 & 0.16 \\
16 & 8 & \textbf{1.00} & 0.15 \\
32 & 4 & 0.68 & 0.16 \\
32 & 8 & \textbf{1.00} & 0.16 \\
64 & 4 & 0.68 & 0.17 \\
64 & 8 & \textbf{1.00} & 0.17 \\
\midrule
\multicolumn{2}{r}{Block-sparse (baseline)} & 0.99 & \textbf{0.77} \\
\bottomrule
\end{tabular}
\end{table}

\Cref{tab:centroid} shows that \textbf{capacity is the dominant variable}:
$\text{cap}=4 \to \text{hit} \approx 0.70$; $\text{cap}=8 \to \text{hit}=1.0$.
Centroid count ($16 \to 64$) has negligible effect. InfoNCE + cap=8 achieves
selector parity with block-sparse (hit$=1.0$), but task accuracy is stuck at
$\sim$0.16 vs.\ block-sparse's 0.77.

\subsection{The Precision/Recall Bottleneck}

With cap=8, each query block selects $c_{pq} \times \text{cap} + 1 = 17$
blocks. The needle is always in the set (hit$=$1.0) but surrounded by 12
extra distractor blocks. We tested three remedies:

\begin{itemize}[leftmargin=*,itemsep=2pt]
\item \textbf{Within-bucket refinement} (top-$k$ from candidates): hit drops
  to 0.72, acc unchanged (0.17).
\item \textbf{Train-high / inference-low} (train cap=8, eval cap=4): hit
  drops to 0.73, acc drops to 0.15.
\item \textbf{Main q/k refinement} (score candidates with attention q/k
  instead of selector): hit 0.72, acc 0.17.
\end{itemize}

None closed the gap. The bottleneck is \emph{not} routing precision but
\emph{representation quality}: the coarse centroid routing path shapes the
attention's q/k/v projections differently than block-sparse's direct pairwise
scoring, limiting token-level extraction even when the needle block is present.

\subsection{LSH Bucketing: Train Dense, Infer Linear}
\label{sec:lsh}

Centroid routing loses information in the \emph{representation} (queries are
compared against lossy cluster summaries). A second linear-time selector, LSH
bucketing~\cite{kitaev2020reformer}, instead prunes the \emph{search}: block representations are hashed by
fixed random projections into buckets ($n_{\text{rounds}}$ rounds), a query
block reads its own bucket, and the candidates are re-scored with the full
$s_q \cdot s_k$ cosine and truncated to block-sparse's budget, so a selected
block is scored at full fidelity. Selection is
$\Ocomplex(n_{\text{blk}} \cdot n_{\text{rounds}} \cdot (n_{\text{buckets}}{+}\text{cap}))$,
linear in $N$; representations are L2-normalized so routing, gate, and hash all
act on the same angular quantity.

Under supervised routing, LSH nearly matches block-sparse and far exceeds
centroid: MQAR accuracy 0.76 vs.\ 0.82 (block-sparse) vs.\ 0.17 (centroid) at
$n_{\text{blk}}{=}16$, and 0.92 vs.\ 0.96 at $n_{\text{blk}}{=}64$, with recall
plateauing at hit$\,{=}\,$0.84 (an irreducible hash-collision miss). Linear
selection \emph{can} therefore preserve accuracy when the representation stays
lossless, which the lossy centroid path cannot. This is the section's main
positive result.

The open part is training without supervision. The routing loss above uses the
known needle block, available only on synthetic MQAR. We tested the unsupervised
regime two ways. On MQAR under gate coupling alone, LSH is fragile: recall is
highly sensitive to the cosine gate temperature (hit 0.09, 0.34, 0.02 at scale
5.66, 30, 50) and never reaches block-sparse's gate-only level (acc 0.96, hit
0.78). On a real language model (TinyStories, 45M, gate coupling only) the result
is sharper (Table~\ref{tab:lsh-lm}): gate-only LSH reaches the \emph{same}
perplexity as own-block-only attention (no routing), and its perplexity does not
improve with context length, whereas block-sparse beats the local baseline by
3.7 perplexity and improves with length. LSH attends to five blocks but adds
nothing over attending to one: its content routing is inert.

\begin{table}[h]
\centering
\caption{Gate-only TinyStories (45M, $B_s{=}16$, $L{=}512$, 3000 steps).
Local-only is top-$k{=}0$ (own block, no routing). LSH matches the no-routing
baseline; only block-sparse routing helps and improves with length.}
\label{tab:lsh-lm}
\small
\begin{tabular}{lcccc}
\toprule
 & val PPL & PPL@128 & PPL@256 & PPL@512 \\
\midrule
Local-only (no routing) & 21.43 & 20.7 & 20.6 & 21.7 \\
LSH (linear routing) & 21.49 & 21.0 & 21.0 & 21.7 \\
Block-sparse routing & \textbf{17.69} & 20.1 & 18.6 & \textbf{17.8} \\
\bottomrule
\end{tabular}
\end{table}

Across these runs one mechanism recurs: hard hashing gives the selector no
gradient on blocks it does not already collide with, so when the selector is
trained \emph{through} the hash it cannot learn to route without supervision,
while block-sparse's all-pairs top-$k$ scores every block on every step and
learns to route in every regime.

This points to what the inference-cost-only objective permits: train through the
dense all-pairs top-$k$ selector (full gradient, identical to block-sparse) and
swap in LSH bucketing only at inference, on the same normalized representations,
so the hash never enters the gradient path.

For \emph{retrieval} this works and scales. On MQAR, LSH inference recovers the
dense-trained selector at a fixed inference budget, and recall holds (even
improves) as context grows: at 8 rounds and 64 buckets (constant, so inference is
$\Ocomplex(n_{\text{blk}})$), $n_{\text{blk}}{=}128/256/512$ give hit
$0.978/0.992/0.994$ (Table~\ref{tab:lsh-scale}). The single needle is reliably
hashed into the query's bucket.

For \emph{language modeling} it does not. At $n_{\text{blk}}{=}128$ (block 4,
$L{=}512$, gate only) and a matched attention budget (own block plus top-4), a
dense selector that scores all blocks reaches 7.9 perplexity, but LSH selecting
its top-4 from $\sim$16 bucket candidates reaches 113 to 1115 (worse at fewer
rounds). The dense selector at the same budget succeeds, so this is not a coverage
floor: LSH's bucketing fails to recall the handful of relevant blocks the dense
selector finds. A single needle is easy to hash to; the blocks a language model
needs per query are not, at least with this hash. The apparent unsupervised LM
recovery held only at near-dense coverage (32 rounds for 32 blocks) and does not
survive at a sparse budget.

So train-dense / infer-LSH yields a genuinely linear, accuracy-preserving selector
for needle-style retrieval, but not yet for general language-model attention. The
open problem is a hash whose recall of content-relevant blocks is high enough for
the LM at a sparse budget; a bit-based SimHash~\cite{rabbani2024lshe} or a learned
hash~\cite{guo2026dashkv} (both used for token-level KV retrieval, the latter
reporting higher recall than plain LSH) are the natural next steps for the
block-level selector.

\begin{table}[h]
\centering
\caption{Train dense, infer LSH at a \emph{fixed} inference budget (8 rounds, 64
buckets, cap 8; $\Ocomplex(n_{\text{blk}})$). Recall holds and improves as context
grows. MQAR, supervised; oracle is the dense-trained selector.}
\label{tab:lsh-scale}
\small
\begin{tabular}{lccc}
\toprule
$n_{\text{blk}}$ & LSH acc & LSH hit & oracle acc \\
\midrule
128 & 0.95 & 0.978 & 0.99 \\
256 & 0.98 & 0.992 & 0.99 \\
512 & 0.99 & 0.994 & 0.99 \\
\bottomrule
\end{tabular}
\end{table}

% ─═ 6. Language Model Experiments ═════════════════════════════════════════
\section{Language Model Experiments}\label{sec:lm}

\subsection{Setup}

\begin{itemize}[leftmargin=*,itemsep=1pt]
\item \textbf{Dataset:} TinyStoriesV2-GPT4~\cite{eldan2023tinystories},
  5M tokens, GPT-2 BPE tokenizer (vocab 50,257).
\item \textbf{Model:} 6 layers, $d=512$, $h=8$, weight-tied embeddings.
  44.9M (dense) / 45.1M (sparse) parameters.
\item \textbf{Training:} 5000 steps, batch 16, $L=512$, AdamW
  ($\beta_1{=}0.9, \beta_2{=}0.95, \text{wd}{=}0.1$), cosine LR $3 \times 10^{-4}$
  with 200-step warmup. bf16 mixed precision.
\item \textbf{Sparse config:} $B_s=64$, top-$k=4$, gate coupling, Triton kernel.
\item \textbf{Hardware:} NVIDIA GB10 Grace-Blackwell, 128\,GB unified memory.
\end{itemize}

\subsection{Results}

\begin{table}[h]
\centering
\caption{TinyStories language model: dense vs.\ block-sparse Triton.}
\label{tab:lm}
\small
\begin{tabular}{lrrrrrr}
\toprule
Model & Params & Val PPL & PPL@128 & PPL@256 & PPL@512 & Time \\
\midrule
Dense & 44.9M & 14.41 & 15.41 & 14.54 & 14.09 & 736\,s \\
Sparse-Triton & 45.1M & 15.00 & 15.31 & \textbf{14.16} & 14.51 & 908\,s \\
\bottomrule
\end{tabular}
\end{table}

The sparse-Triton model achieves val perplexity 15.0 versus 14.4 for
dense, a 4\% gap. Notably, at 256-token context, sparse \emph{beats} dense
(14.16 vs.\ 14.54), suggesting the learned selector provides a useful
inductive bias at moderate lengths. The training time overhead (908\,s vs.\
736\,s, $+23\%$) is from the selector computation and the Triton kernel's
contiguity copies.

\subsection{NIAH (Needle-in-a-Haystack)}

Both models score 0\% on NIAH at $L=512$ (50 samples). This is a model
capacity limitation (45M params, 3000 to 5000 steps), not a sparse attention
defect. NIAH retrieval requires 100M+ parameters and 10k+ training steps.

% ─═ 7. Benchmark ══════════════════════════════════════════════════════════
\section{Benchmark: Latency and Memory}\label{sec:bench}

\begin{table}[h]
\centering
\caption{Forward-pass benchmark: $D{=}512$, $H{=}8$, $B{=}1$, $B_s{=}128$,
$k{=}8$, causal. GB10 GPU.}
\label{tab:bench}
\small
\begin{tabular}{rrrrrrr}
\toprule
$L$ & Dense (ms) & Sparse-PT (ms) & Triton (ms) & Dense (GB) & PT (GB) & Triton (GB) \\
\midrule
4,096 & 0.25 & 11.64 & 1.57 & 0.02 & 0.70 & 0.05 \\
16,384 & 3.05 & 44.08 & 6.80 & 0.10 & 2.70 & 0.10 \\
65,536 & 48.20 & 178.59 & 28.32 & 0.30 & 10.70 & 0.30 \\
131,072 & 195.38 & 365.72 & 56.79 & 0.58 & 21.38 & 0.57 \\
262,144 & 796.59 & 739.32 & \textbf{114.37} & 1.12 & 42.72 & \textbf{1.11} \\
\bottomrule
\end{tabular}
\end{table}

\paragraph{Memory.} The Triton kernel uses 1.11\,GB at 262k -parity with
dense FlashAttention (1.12\,GB) and $40\times$ less than the PyTorch
gather reference (42.72\,GB). The fused kernel never materializes the
gathered K/V tensor; all intermediate state lives in SRAM.

\paragraph{Latency.} The Triton kernel is $7.4\times$ faster than dense at
262k (114\,ms vs.\ 797\,ms). The crossover occurs at $\sim$32k tokens. At
65k, the kernel is already $1.8\times$ faster than dense.

\paragraph{PyTorch gather is the bottleneck.} The same algorithm, same
selection, but PyTorch materializes the gather (42.7\,GB) and is $6.5\times$
slower than Triton at 262k. The fused kernel proves the sparse mechanism
itself is strictly superior to dense at long contexts; the reference
implementation was the obstacle.

\subsection{Selection vs.\ the Attention Read: the $N^{4/3}$ Floor}
\label{sec:floor}

The benchmark above measures the fused attention \emph{read}, which is linear.
The block-pair \emph{selection} (scoring all $n_{\text{blk}}^2$ block pairs,
\Cref{sec:method}) is not. We isolate the two stages and sweep $L$ to 16M
tokens across block sizes (\Cref{tab:floor}).

\begin{table}[h]
\centering
\caption{Selection (quadratic) vs.\ attention read (linear), isolated.
$D{=}512$, $H{=}8$, $B_s{=}64$, $k{=}8$, causal, GB10 (130\,GB). ``OOM'': the
$\Ocomplex(n_{\text{blk}}^2)$ score matrix exceeds memory.}
\label{tab:floor}
\small
\begin{tabular}{rrrr}
\toprule
$L$ & Selection (ms) & Read (ms) & Select mem (GB) \\
\midrule
262,144 & 5.8 & 17.1 & 0.4 \\
1,048,576 & 51.2 & 84.7 & 3.0 \\
4,194,304 & 652.3 & 359.1 & 34.4 \\
8,388,608 & \textbf{OOM} & 716.0 & $>$130 \\
16,777,216 & \textbf{OOM} & 1{,}447.6 & $>$130 \\
\bottomrule
\end{tabular}
\end{table}

\paragraph{Scaling.} Fitting across $L$, the attention read scales with
exponent $0.96$ to $1.02$ (linear) at every block size, while selection scales
with exponent $1.5$ to $1.9$ (quadratic, approaching $2$ as $N$ leaves the
latency floor). Selection latency overtakes the read at $\sim$512k tokens for
$B_s{=}32$ and $\sim$2M for $B_s{=}64$; for $B_s{=}128$ it stays below the read
through 8M (ratio $0.17$), crossing over only in the hundreds of millions.

\paragraph{The quadratic dies by OOM, not by latency.} The
$\Ocomplex(n_{\text{blk}}^2)$ score matrix is a memory wall reached well before
the latency crossover: it OOMs at 4M ($B_s{=}32$), 8M ($B_s{=}64$), and 16M
($B_s{=}128$), while the linear read runs to 16M (86\,GB) on the same GPU. In
practice, quadratic selection becomes impossible to hold, not merely slow.

\paragraph{Variable block size does not rescue linearity.} A natural trick is
to grow the block with context, $B_s \propto \sqrt{N}$, making selection
$\Ocomplex((N/\sqrt{N})^2) = \Ocomplex(N)$, linear. We confirm this: along a
$\sqrt{N}$ schedule ($B_s{=}32,64,128$ at $L{=}$256k, 1M, 4M) the measured
selection exponent falls to $0.98$. But the read processes $k B_s$ keys per
query, so it then grows as $N B_s \propto N^{1.5}$ (measured exponent rising to
$2.3$): the quadratic merely relocates from selection into the read. Writing
the total as $N^2 s / B_s^2 + N c B_s$ and minimizing over $B_s$ gives
$B_s \propto N^{1/3}$, at which \emph{both} terms equal $\Ocomplex(N^{4/3})$.
This $N^{4/3}$ is a floor: \textbf{no block-size schedule makes single-level,
all-pairs block selection linear end-to-end; the quadratic can be moved between
stages but not removed.} Genuinely linear selection requires scoring
$\Ocomplex(C)$ candidates per query (clustering, hashing, or hierarchy), not
all $n_{\text{blk}}$, the open problem of \Cref{sec:centroid}.

% ─═ 8. Discussion ═════════════════════════════════════════════════════════
\section{Discussion}\label{sec:discussion}

\subsection{What Works}

Gate coupling + auxiliary routing loss reliably trains the block selector
(hit $0.99$ on MQAR, val PPL within 4\% of dense on TinyStories
\cite{eldan2023tinystories}). The fused Triton kernel delivers the theoretical
memory and latency advantages of block-sparse attention in practice, with no
accuracy loss from numerical approximation.

\subsection{What Doesn't: Centroid Routing}

Centroid routing achieves linear selection cost and perfect retrieval recall
(hit$=$1.0) but cannot match block-sparse task accuracy (acc 0.16 vs.\ 0.77).
The failure is not in routing (InfoNCE \cite{oord2018infonce} fixes that) but
in \emph{representation quality}: the coarse centroid path shapes the
attention projections differently, limiting fine-grained token extraction.
This is a fundamental limitation of content-routed linear-selection attention
that, to our knowledge, has not been previously characterized.

\subsection{Limitations}

\begin{itemize}[leftmargin=*,itemsep=2pt]
\item Single-level block-pair selection is $\Ocomplex(n_{\text{blk}}^2)$ with
  an $\Ocomplex(N^{4/3})$ floor no block schedule can beat (\Cref{sec:floor});
  it exhausts memory before any latency crossover. Centroid routing is linear
  but accuracy-broken.
\item NIAH requires a larger model; our 45M model is too small.
\item The backward kernel uses \texttt{atomic\_add} for dK/dV, which may
  serialize under high contention. A split-K or warp-level reduction could
  improve throughput.
\item No autoregressive KV-cache support; the kernel is training/prefill only.
\end{itemize}

% ─═ 9. Conclusion ═════════════════════════════════════════════════════════
\section{Conclusion}\label{sec:conclusion}

We demonstrated that block-sparse attention with a fused Triton kernel
achieves memory parity with dense FlashAttention and $7.4\times$ latency
speedup at long contexts, while maintaining within-4\% perplexity on a real
language model. The gate-coupled selector training recipe is simple,
effective, and requires no separate pre-training stage.

We probed two linear-time selectors. Centroid routing reaches linear cost and
perfect recall but degraded accuracy: the obstacle is representation quality
under coarse routing, not routing recall (InfoNCE resolves the latter). LSH
bucketing, which prunes the search rather than the representation, matches
block-sparse accuracy under supervised routing and far exceeds centroid, showing
a lossless-refine linear selector can preserve accuracy; but its unsupervised
routing proved inert, matching a no-routing baseline on a 45M language model.
Reliable unsupervised linear selection remains open.

Promising directions include scaling to the 100M+ regime where
needle-in-a-haystack retrieval becomes measurable, extending the kernel to
autoregressive KV-cache decoding, and designing a linear-selection mechanism
that preserves the representation quality of pairwise scoring. More broadly,
the fused kernel removes the implementation barrier that has masked
block-sparse attention's asymptotic advantage: at long contexts,
content-dependent sparsity is now both faster and lighter than dense
attention in practice, not merely in theory.

% ─═ References ════════════════════════════════════════════════════════════
\bibliographystyle{plain}
\begin{thebibliography}{99}

\bibitem{dao2022flashattention}
T.~Dao, D.~Y.~Fu, S.~Ermon, A.~Rudra, and C.~R\'e.
\newblock {FlashAttention}: Fast and memory-efficient exact attention with
  {IO}-awareness.
\newblock In \emph{Advances in Neural Information Processing Systems
  (NeurIPS)}, 2022. arXiv:2205.14135.

\bibitem{dao2023flashattention2}
T.~Dao.
\newblock {FlashAttention-2}: Faster attention with better parallelism and work
  partitioning.
\newblock In \emph{International Conference on Learning Representations
  (ICLR)}, 2024. arXiv:2307.08691.

\bibitem{yuan2025nsa}
J.~Yuan, H.~Gao, D.~Dai, J.~Luo, L.~Zhao, Z.~Zhang, Z.~Xie, Y.-X.~Wei,
  L.~Wang, Z.~Xiao, Y.~Wang, C.~Ruan, M.~Zhang, W.~Liang, and W.~Zeng.
\newblock Native sparse attention: Hardware-aligned and natively trainable
  sparse attention.
\newblock \emph{arXiv preprint} arXiv:2502.11089, 2025. ACL 2025 Best Paper.

\bibitem{lu2025moba}
E.~Lu, Y.~Duan, H.~Hu, W.~Wang, J.~Zhu, H.~Liu, K.~Sun, and T.~Wang.
\newblock {MoBA}: Mixture of block attention for long-context {LLMs}.
\newblock In \emph{NeurIPS}, 2025. arXiv:2502.13189.

\bibitem{roy2021routing}
A.~Roy, M.~Saffar, A.~Vaswani, and D.~Grangier.
\newblock Efficient content-based sparse attention with routing transforms.
\newblock \emph{Transactions of the Association for Computational Linguistics
  (TACL)}, 9:53-68, 2021.

\bibitem{shazeer2017moe}
N.~Shazeer, A.~Mirhoseini, K.~Maziarz, A.~Davis, Q.~V.~Le, G.~E.~Hinton, and
  J.~Dean.
\newblock Outrageously large neural networks: The sparsely-gated
  mixture-of-experts layer.
\newblock In \emph{International Conference on Learning Representations
  (ICLR)}, 2017. arXiv:1701.06538.

\bibitem{beltagy2020longformer}
I.~Beltagy, M.~E.~Peters, and A.~Cohan.
\newblock {Longformer}: The long-document transformer.
\newblock \emph{arXiv preprint} arXiv:2004.05150, 2020.

\bibitem{zaheer2020bigbird}
M.~Zaheer, G.~Guruganesh, A.~Dubey, J.~Ainslie, C.~Alberti, S.~Onta\~{n}\'on,
  et~al.
\newblock {BigBird}: Transformers for longer sequences.
\newblock In \emph{NeurIPS}, 2020. arXiv:2007.14062.

\bibitem{kitaev2020reformer}
N.~Kitaev, \L.~Kaiser, and A.~Levskaya.
\newblock {Reformer}: The efficient transformer.
\newblock In \emph{ICLR}, 2020. arXiv:2001.04451.

\bibitem{tay2020sinkhorn}
Y.~Tay, D.~Bahri, D.~Metzler, D.-C.~Juan, Z.~Zhao, and C.~Zheng.
\newblock Sparse sinkhorn attention.
\newblock In \emph{ICML}, 2020. arXiv:2002.11296.

\bibitem{olsson2022induction}
C.~Olsson, N.~Elhage, N.~Nanda, N.~Joseph, N.~DasSarma, T.~Henighan, B.~Mann,
  et~al.
\newblock In-context learning and induction heads.
\newblock \emph{Transformer Circuits Thread}, Anthropic, 2022.

\bibitem{eldan2023tinystories}
R.~Eldan and Y.~Li.
\newblock {TinyStories}: How small can language models be and still speak
  coherent {English}?
\newblock \emph{arXiv preprint} arXiv:2305.07759, 2023.

\bibitem{tillet2019triton}
P.~Tillet, H.-T.~Kung, and D.~D.~Cox.
\newblock {Triton}: An intermediate language and compiler for tiled neural
  network computations.
\newblock In \emph{Proc.\ 3rd ACM SIGPLAN Int.\ Workshop on Machine Learning
  and Programming Languages (MAPL)}, pages 10-19, 2019.

\bibitem{oord2018infonce}
A.~van den Oord, Y.~Li, and O.~Vinyals.
\newblock Representation learning with contrastive predictive coding.
\newblock \emph{arXiv preprint} arXiv:1807.03748, 2018.

\bibitem{arora2024mqar}
S.~Arora, S.~Gopi, N.~Ingster, and P.~Nakkiran.
\newblock Simple, fast, and practically good: A simple {LLM}-based sparse
  retrieval model for {MQAR}.
\newblock 2024.

\bibitem{tang2024quest}
J.~Tang, Y.~Zhao, K.~Zhu, G.~Xiao, B.~Kasikci, and S.~Han.
\newblock {Quest}: Query-aware sparsity for efficient long-context {LLM}
  inference.
\newblock In \emph{ICML}, 2024. arXiv:2406.10774.

\bibitem{zhang2023h2o}
Z.~Zhang, Y.~Sheng, T.~Zhou, T.~Chen, L.~Zheng, R.~Cai, Z.~Song, Y.~Tian,
  C.~R\'e, C.~Barrett, Z.~Wang, and B.~Chen.
\newblock {H2O}: Heavy-hitter oracle for efficient generative inference of large
  language models.
\newblock In \emph{NeurIPS}, 2023. arXiv:2306.14048.

\bibitem{liu2024retrievalattention}
D.~Liu, M.~Chen, B.~Lu, H.~Jiang, Z.~Han, Q.~Zhang, Q.~Chen, C.~Zhang, B.~Ding,
  K.~Zhang, C.~Chen, F.~Yang, Y.~Yang, and L.~Qiu.
\newblock {RetrievalAttention}: Accelerating long-context {LLM} inference via
  vector retrieval.
\newblock \emph{arXiv preprint} arXiv:2409.10516, 2024.

\bibitem{rabbani2024lshe}
T.~Rabbani, M.~Liu, T.~O'Halloran, A.~Sankaralingam, M.-A.~Hartley, and
  F.~Huang.
\newblock {LSH-E}: An adaptive locality-sensitive strategy for {KV} cache
  compression.
\newblock In \emph{NeurIPS}, 2024.

\bibitem{guo2026dashkv}
J.~Guo, Z.~Zhang, J.~Xie, M.~T.~Iqbal, D.~Han, L.-H.~Lee, S.-H.~Bae, J.~Zou,
  Y.~Yang, and C.~Zhang.
\newblock {DASH-KV}: Accelerating long-context {LLM} inference via asymmetric
  {KV} cache hashing.
\newblock \emph{arXiv preprint} arXiv:2604.19351, 2026.

\end{thebibliography}

\end{document}
